\documentclass[11pt]{article}
\usepackage{amsmath,amssymb}
\usepackage{physics}
\usepackage[margin=1in]{geometry}
\usepackage{enumitem}
\usepackage{listings}

\lstset{basicstyle=\ttfamily\small, columns=fullflexible, keywordstyle=\bfseries}

\title{ATPL 2025 --- Exercises: Quantum Programs}
\author{Introduction to Quantum Computing}
\date{}

\begin{document}
\maketitle

Background: the \texttt{Step}/\texttt{Program} model and \texttt{Op} data
type from the lecture, and Nielsen \& Chuang Section 4.2--4.3 on
universal gate sets.

\begin{enumerate}[label=\textbf{Exercise \arabic*.}, leftmargin=2.2cm]

\item \textbf{Unitarity from probability conservation.}
Redo the argument from the lecture in your own words: given a pure
quantum program $P$ acting on an $n$-qubit state, explain precisely why
requiring $\sum_i |a_i|^2 = \sum_i |b_i|^2 = 1$ for \emph{every} possible
input state $\ket\Psi=\sum_i a_i\ket i$ forces $P$ (as a matrix) to
satisfy $P^\dagger P = I$. (Hint: consider
$\braket{\Psi}{\Psi} = \bra\Psi P^\dagger P\ket\Psi$ and use that this
must equal $1$ for all normalized $\ket\Psi$.)

\item \textbf{Reading circuits as \texttt{Op} expressions.}
Using the grammar
\[
\begin{aligned}
\texttt{Op} ::= {}&
    I \mid X \mid SX \mid Y \mid Z \mid H \mid R\ \texttt{Op}\ \theta \\
  &\mid C\ \texttt{Op} \mid SWAP \\
  &\mid \texttt{Tensor}\ \texttt{Op}\ \texttt{Op} \\
  &\mid \texttt{DirectSum}\ \texttt{Op}\ \texttt{Op} \\
  &\mid \texttt{Compose}\ \texttt{Op}\ \texttt{Op}.
\end{aligned}
\]
write down the \texttt{Op} expression, and the corresponding matrix, for:
\begin{enumerate}[label=(\alph*)]
\item Applying $H$ to qubit 1 and doing nothing to qubit 2
(2-qubit operator).
\item CNOT with qubit 1 as control and qubit 2 as target, expressed using
\texttt{C}.
\item The circuit that prepares the Bell state $\ket{\Phi^+}$ from
$\ket{00}$ (Hadamard on qubit 1, then CNOT).
\end{enumerate}

\item \textbf{Programs as sequences of Steps.}
Using the \texttt{data Step = Op | Measure Int} and
\texttt{data Program = List Step} definitions from the lecture, write
out (as a list of \texttt{Step}s, in pseudo-Haskell) a 2-qubit program
that:
\begin{enumerate}[label=(\alph*)]
\item Prepares $\ket{\Phi^+}$ and then measures qubit 1.
\item Prepares $\ket{\Phi^+}$, measures qubit 1, and then \emph{also}
measures qubit 2. What outcome distribution do you expect for the second
measurement, given each possible outcome of the first? Justify using the
measurement exercises.
\end{enumerate}

\item \textbf{DirectSum and multiplexed operations.}
The slide describes \texttt{DirectSum} as "multiplexed operations —
generalizing conditional execution." Consider
$U = \texttt{DirectSum}\ U_0\ U_1$ where $U_0,U_1$ are 1-qubit unitaries,
acting on a 2-qubit system where the first qubit is the "control."
\begin{enumerate}[label=(\alph*)]
\item Write $U$ as a $4\times4$ block-diagonal matrix in terms of the
$2\times2$ matrices for $U_0$ and $U_1$.
\item Show that $\texttt{DirectSum}\ I\ X$ is exactly CNOT (with the
convention from Exercise 4 above). What is $\texttt{DirectSum}\ I\ I$?
\item Explain in one or two sentences why \texttt{DirectSum} is a
natural generalization of classical \texttt{if/else} branching to the
quantum setting, and why it must still be built from a \emph{unitary}
$U_0$ and $U_1$ rather than arbitrary classical branches.
\end{enumerate}

\item \textbf{Universality (conceptual, tie-together).}
The lecture claims that one general single-qubit rotation, one fixed
orthogonal rotation by $\pi/2$, and CNOT are enough to build any pure
quantum program.
\begin{enumerate}[label=(\alph*)]
\item Using only $\{H, R_Z(\theta) \text{ for any } \theta, \mathrm{CNOT}\}$,
sketch how you would build an arbitrary single-qubit unitary $U$ (recall
the $\mathbf n\cdot\boldsymbol\sigma$ decomposition from the Pauli/Bloch
exercises, and that $H$ conjugates $Z$-rotations into $X$-rotations).
\item Why is a 2-qubit entangling gate like CNOT specifically necessary
--- i.e.\ why can no sequence of \texttt{Tensor}-only (single-qubit)
operations ever prepare an entangled state such as $\ket{\Phi^+}$? (You
already proved the key fact in the tensor-product exercises.)
\end{enumerate}

\end{enumerate}

\end{document}
